\chapter{Related Work and State of the Art}
\label{chap:related}

In this chapter we describe related work in the area of protocol fuzzing and the current state of the art.

The field of protocol fuzzing has evolved significantly, with various techniques aimed at uncovering vulnerabilities in network communications. A comprehensive survey by Chen et al.~\cite{chen2024survey} categorizes protocol fuzzing methods into generation-based, mutation-based, and hybrid approaches, highlighting the role of state-aware fuzzing in handling complex protocol interactions. The authors discuss challenges such as state explosion and the need for efficient seed selection, noting that while traditional black-box fuzzers like Peach and Sulley were foundational, modern tools incorporate AFL-inspired feedback mechanisms to improve coverage. They emphasize that protocol fuzzing differs from general software fuzzing due to the need for modeling message formats and sequences, often requiring domain-specific knowledge to avoid invalid test cases that are immediately rejected. The survey also covers evaluation metrics such as bug discovery rates and coverage depth, and points to limitations in handling encrypted protocols like QUIC, which underlies WebTransport, where decryption hooks or proxying are needed.

Building on this, a study by Wang et al.~\cite{wang2024enhancing} focuses on fuzzing for network security, proposing a framework that integrates machine learning for smarter mutation strategies. They distinguish black-box approaches, which rely solely on input-output observations, from gray-box (with partial knowledge) and white-box (full instrumentation), arguing that gray-box achieves a balance for protocols where source code access is limited. In their experiments testing against HTTP/2 and TLS implementations, they found that sequence duplication and frame reordering revealed more state-related bugs than random mutations. They note, however, gaps in fuzzing multiplexed protocols, where concurrent streams complicate state tracking, which is directly relevant to the design of WebTransport. The work criticizes existing tools for lacking support for emerging IETF standards, suggesting that custom extensions are often needed, as was done in this thesis with BooFuzz.

Daniele et al.~\cite{daniele2024libafl} present LibAFL*, an advanced fuzzing library that supports custom harnesses for protocol testing, demonstrating its use on industrial control systems. They emphasize the importance of non-invasive monitoring, such as log analysis, over invasive instrumentation, which aligns with our language-agnostic approach. The paper compares LibAFL* with BooFuzz, noting BooFuzz's strengths in ease of use for Python-based scripting but weaknesses in handling high-speed protocols without asynchronous integrations. Their findings show that without instrumentation, crash detection relies on external signals such as process termination or log errors, which we adopted for WebTransport servers. A key gap identified is the lack of fuzzers for web-oriented protocols beyond HTTP, with no mention of WebTransport, which underscores the novelty of our contribution.

Finally, a recent preprint by Li et al.~\cite{li2024fuzzing} surveys fuzzing techniques for IoT and network protocols, highlighting the rise of model-based fuzzing, where protocol specifications (e.g., from IETF drafts) guide test case generation. They discuss tools such as Fuzzowski and BooFuzz, praising BooFuzz for its block-based message modeling but criticizing its lack of built-in support for sequence mutations such as duplication, which required modifications in our implementation. The authors point out that while fuzzing has been extensively applied to mature protocols such as TCP/IP or MQTT, emerging ones like WebTransport---still in draft stage but implemented in browsers---remain understudied, with potential vulnerabilities in QUIC integration untested. They suggest community-driven tools on GitHub, similar to our project, to fill these gaps.

\section{Comparison with Our Work}

Compared to these works, our thesis addresses a specific gap in WebTransport fuzzing that none of the surveyed papers directly address, likely due to the protocol's novelty. While surveys such as~\cite{chen2024survey} and~\cite{li2024fuzzing} provide broad overviews, they do not include WebTransport examples, and tools like BooFuzz require extensions for QUIC-based protocols, as we implemented using aioquic. Unlike the white-box approaches in~\cite{daniele2024libafl}, our black/gray-box method avoids language-specific instrumentation, making it more universal but potentially less deep in coverage; however, it proved effective in discovering issues such as a reachable assertion bug in the Rust wtransport library~\cite{wtransport_github}, where sending a drain capsule without a proper HTTP/3 frame caused a worker thread panic. This contrasts with related fuzzing studies on HTTP/3 that focus on the client side but overlook server echo implementations.

Overall, the state of the art reveals a mature ecosystem for protocol fuzzing but with limited application to new web transports, which motivated our development of a BooFuzz-based fuzzer tailored for WebTransport message structures and sequences.
